\subsection{Task Formulation}
\label{sec:task}

Formally, given an input sentence $s = w_1, w_2, \ldots, w_n$, AMR parsing
requires predicting a directed graph $G = (V, E, L_v, L_e)$ where $V$ is a
set of concept nodes, $E$ is a set of directed edges, $L_v : V \to \mathcal{C}$
assigns concept labels from an open vocabulary $\mathcal{C}$, and
$L_e : E \to \mathcal{R}$ assigns relation labels from a closed set $\mathcal{R}$
(e.g., \texttt{:ARG0}, \texttt{:ARG1}, \texttt{:mod}). Variable identifiers
(e.g., $l$, $b$) are auxiliary; their specific values do not affect graph
semantics, but they are necessary to encode reentrancy---the case where a single
node participates as the argument of multiple relations.

\begin{figure}[h]
\begin{lstlisting}[style=amrstyle, basicstyle=\small\ttfamily]
Input: "The boy left."

AMR: (l / leave-01
        :ARG0 (b / boy))
\end{lstlisting}
\caption{A minimal AMR graph. The root concept \texttt{leave-01} is linked
to \texttt{boy} via the \texttt{:ARG0} (agent) relation. Variables $l$ and
$b$ are arbitrary identifiers; reentrancy is encoded by reusing a variable in
multiple edge positions.}
\label{fig:amr-example}
\end{figure}

End-to-end parsers produce $G$ directly from $s$ in a single decoding pass,
conflating concept identification, relation labelling, and variable assignment.
We argue that this joint formulation imposes an unnecessarily high structural
load on the decoder, and propose to factor it into explicit intermediate
representations.

\subsection{Multi-Task Unified Prompting (MTUP)}
\label{sec:mtup}

\paragraph{Overview.}
MTUP is a supervised fine-tuning framework in which a single instruction-tuned
language model is trained to produce a \emph{sequence of intermediate
representations} culminating in the final AMR graph, all within one assistant
turn. Decomposition is implemented through structured step markers embedded in
the target text (e.g., \texttt{[STEP 1: AMR\_NO\_VAR]},
\texttt{[STEP 2: AMR\_WITH\_VAR]}). At inference time, the model generates
all steps autoregressively; only the final step is extracted for evaluation.

\paragraph{Training objective.}
Each training example is formatted as a multi-turn chat sequence using
Qwen's \texttt{<|im\_start|>...<|im\_end|>} template. The system turn provides
task instructions and language-specific AMR conventions; the user turn contains
the input sentence; the assistant turn contains the multi-step target. Loss is
computed exclusively over assistant tokens, masking both the system and user
turns. This ensures the model learns to generate structured semantic outputs
rather than to reproduce instructions.

\paragraph{Advantage of unified prompting.}
Training a single model over all steps---rather than a pipeline of separate
models---allows the model to learn cross-step dependencies during training, and
requires only a single adapter at inference time. The same training codebase
and model architecture applies without modification to Vietnamese, English, and
any future target language.

\subsection{MTUP Variants}
\label{sec:variants}

We evaluate four decomposition configurations. All share the same base model
and LoRA configuration; they differ only in the number and type of intermediate
steps.

\subsubsection{MTUP2: Two-Step Decomposition}

\begin{lstlisting}[style=amrstyle, basicstyle=\scriptsize\ttfamily]
[STEP 1: AMR_NO_VAR]
(leave-01 :ARG0 (boy))

[STEP 2: AMR_WITH_VAR]
(l / leave-01 :ARG0 (b / boy))
\end{lstlisting}

MTUP2 separates the \emph{structural} sub-problem---predicting the set of
concept nodes and their relational skeleton---from the \emph{binding}
sub-problem of assigning variable identifiers. In the no-variable form,
reentrancy is represented by repeating the concept label at each occurrence,
which eliminates the requirement to track global variable scope during the
first generation step. We hypothesise that this factorisation reduces the
probability of syntactic errors (unbalanced parentheses, undefined variables)
by deferring their most frequent source to the final step.

Among all tested variants, MTUP2 achieves the best SMATCH score on Vietnamese
(0.6494 mean over 3 seeds). We therefore treat it as our primary proposed
method, while acknowledging that the advantage over B2 (0.6466) is small and
not statistically established under our evaluation protocol.

\subsubsection{MTUP3A: Three Steps with Implicit Relations}

\begin{lstlisting}[style=amrstyle, basicstyle=\scriptsize\ttfamily]
[STEP 1: CONCEPTS]
concept_1 | leave-01
concept_2 | boy

[STEP 2: AMR_NO_VAR]
(leave-01 :ARG0 (boy))

[STEP 3: AMR_WITH_VAR]
(l / leave-01 :ARG0 (b / boy))
\end{lstlisting}

MTUP3A adds an explicit concept enumeration step before graph construction.
By committing to the set of semantic elements prior to relation assignment, we
hypothesised that the model could avoid generating spurious concepts. Empirically,
MTUP3A achieves SMATCH 0.6354 on Vietnamese---below MTUP2 by 1.4pp---suggesting
that the concept-listing step introduces noise or over-constrains the subsequent
generation.

\subsubsection{MTUP3B: Three Steps with Explicit Relations}

\begin{lstlisting}[style=amrstyle, basicstyle=\scriptsize\ttfamily]
[STEP 1: CONCEPTS]
concept_1 | leave-01
concept_2 | boy

[STEP 2: RELATIONS]
leave-01 | :ARG0 | boy

[STEP 3: AMR_WITH_VAR]
(l / leave-01 :ARG0 (b / boy))
\end{lstlisting}

MTUP3B makes both concepts and relations explicit before final graph
generation, replacing the no-variable AMR intermediate with a flat relation
list. The hypothesis was that an explicit relation-prediction step would focus
model attention on edge labels and reduce argument attachment errors.
MTUP3B achieves 0.5891 SMATCH on Vietnamese, the lowest among all MTUP
variants, performing similarly to the weaker B1 baseline. We attribute this to
the flat relation list format being a novel output form not represented in
pre-training data, increasing the effective learning burden.

\subsubsection{MTUP4: Four-Step Full Decomposition}

\begin{lstlisting}[style=amrstyle, basicstyle=\scriptsize\ttfamily]
[STEP 1: CONCEPTS]
concept_1 | leave-01
concept_2 | boy

[STEP 2: RELATIONS]
leave-01 :ARG0 boy

[STEP 3: AMR_NO_VAR]
(leave-01 :ARG0 (boy))

[STEP 4: AMR_WITH_VAR]
(l / leave-01 :ARG0 (b / boy))
\end{lstlisting}

MTUP4 combines explicit concept and relation enumeration (as in MTUP3B) with
an intermediate graph reconstruction step (as in MTUP2). It achieves SMATCH
0.5970 on Vietnamese---better than MTUP3B but below MTUP2---suggesting that
the additional Penman intermediate step partially mitigates the errors
introduced by explicit list steps, without recovering to the MTUP2 level.

\subsection{Baselines}
\label{sec:baselines}

\paragraph{B1: Direct generation.}
The model generates the full AMR with variable bindings in a single step,
prefixed by the token \texttt{AMR:}. This represents the standard end-to-end
decoding paradigm and provides the lower reference point for decomposition.

\paragraph{B2: No-variable generation with post-processing.}
The model generates AMR without variable bindings (\texttt{AMR\_NO\_VAR:}),
and a deterministic post-processing routine greedily assigns variables by
first-mention order. This isolates the structural generation quality from
variable assignment and serves as an ablation of MTUP2's first step.

\subsection{Implementation}
\label{sec:implementation}

\paragraph{Model.}
All experiments use Qwen2.5-7B-Instruct with LoRA \citep{hu2021lora}:
rank $r=128$, $\alpha=256$, dropout 0.05, applied to all attention and
feed-forward projection matrices. Flash Attention 2 is enabled for GPU memory
efficiency.

\paragraph{Training.}
We use AdamW with a peak learning rate of $5 \times 10^{-5}$ and cosine decay,
warmup ratio 0.05, effective batch size 16 (1 example per device,
16 gradient accumulation steps), and 25 epochs. Mixed precision uses
\texttt{bf16} where available, with \texttt{fp16} fallback. All experiments
are replicated with three random seeds (42, 123, 456). Loss is masked to
assistant tokens only.

\textbf{Caveat on epoch count.} Training for 25 epochs on $\sim$2{,}950
examples is aggressive relative to standard convention for datasets of this
size. Vietnamese experiments lack a held-out development set for early stopping,
so convergence cannot be monitored during training. Results should be
interpreted in light of this potential for overfitting, and future work should
incorporate a development split with early stopping.

\paragraph{Inference.}
At inference time, each test sentence is formatted as a user turn with the
same system prompt used during training. The model generates autoregressively
up to a method-specific token budget (B1/B2: 512; MTUP2/MTUP3A/MTUP3B: 1536;
MTUP4: 2048). A method-specific extraction routine identifies the final AMR
step in the output; malformed or incomplete outputs are replaced with the
placeholder \texttt{(a / amr-unknown)}.
